\section*{Chapter \ref{ch:g7:volcanoes} --- Volcanoes}

\begin{solution}{exo:g7:volcanoes:1}
(a) \omterm{def:g7:volcanoes:vent}{crater}; (b) \omterm{def:g7:volcanoes:magma}{magma}; (c) \omterm{def:g7:volcanoes:vent}{vent}; (d) \omterm{def:g4:volcanoes-earthquakes-first-look:eruption}{lava}; (e) \omterm{def:g7:volcanoes:magma-chamber}{magma chamber}.
\end{solution}

\begin{solution}{exo:g7:volcanoes:2}
(a) Effusive. (b) Explosive. (c) Effusive: a \omterm{def:g4:volcanoes-earthquakes-first-look:eruption}{lava} fountain throws out
glowing drops which fall back and flow; the \omterm{def:g7:volcanoes:magma}{magma} is not blown into ash.
(d) Explosive: it is a \omterm{def:g7:volcanoes:pyroclastic-flow}{pyroclastic flow}.
\end{solution}

\begin{solution}{exo:g7:volcanoes:3}
(a) True. (b) False: the most abundant volcanic gas is water vapour;
carbon dioxide comes next. (c) False: a \omterm{def:g7:volcanoes:pyroclastic-flow}{pyroclastic flow} is an avalanche
of hot gas, ash and pieces of \omterm{def:g3:stones-pebbles-sand:rock}{rock}. (d) False: they stand inside the
Pacific \omterm{def:g7:plate-tectonics:plate}{plate}, more than 3200~km from the nearest boundary, above a
\omterm{def:g7:volcanoes:hotspot}{hotspot}.
\end{solution}

\begin{solution}{exo:g7:volcanoes:4}
(a) About 1000 to 1200~degrees. (b) About 600 to 800~degrees. (c) Near the
\omterm{def:g7:volcanoes:vent}{vent}, where the \omterm{def:g4:volcanoes-earthquakes-first-look:eruption}{lava} is yellow (1000 to 1200~degrees) rather than orange
(800 to 1000~degrees): the \omterm{def:g4:volcanoes-earthquakes-first-look:eruption}{lava} cools as it flows away from the \omterm{def:g7:volcanoes:vent}{vent}.
\end{solution}

\begin{solution}{exo:g7:volcanoes:5}
(a) $838 + 221 = 1059$ \omterm{def:g4:volcanoes-earthquakes-first-look:volcano}{volcanoes}. (b) $1059 \div 1214 \approx 0.872$:
about \textbf{87~\%}. (c) $69 + 18 + 12 + 0.5 = 99.5$, about 100~\%: the
small difference comes from rounding each percentage.
\end{solution}

\begin{solution}{exo:g7:volcanoes:6}
(a) $564 \div 660 \approx 0.855$: about \textbf{85~\%}. (b)
$660 - 564 - 62 - 33 = 1$: one steep cone. (c) Steep cones are built by
sticky \omterm{def:g4:volcanoes-earthquakes-first-look:eruption}{lava} and ash, the mark of \omterm{def:g7:volcanoes:eruption-types}{explosive eruptions}: along the Ring of
Fire, where one \omterm{def:g7:plate-tectonics:plate}{plate} sinks under another, we expect many \omterm{def:g7:volcanoes:eruption-types}{explosive
eruptions}, like those of Mount St.~Helens and Pinatubo.
\end{solution}

\begin{solution}{exo:g7:volcanoes:7}
(a) $1170 - 1140 = 30$~degrees. (b) $30 \div 10 = 3$~degrees per
kilometre. (c) In the tunnel the \omterm{def:g4:volcanoes-earthquakes-first-look:eruption}{lava} does not touch the air: the
hardened roof above it keeps the heat in, like the lid of a pot. In the
open, its surface loses heat to the air at once.
\end{solution}

\begin{solution}{exo:g7:volcanoes:8}
(a) Kauai, by at least $5.5 - 0.7 = 4.8$~million years. (b) Kauai: \omterm{def:g4:clouds-rain-snow:precipitation}{rain}
and waves have been wearing it down for much longer. (c) The Big Island,
where new \omterm{def:g4:volcanoes-earthquakes-first-look:eruption}{lava} still comes out. (d) Towards the north-west: Kauai, carried
away from the \omterm{def:g7:volcanoes:hotspot}{hotspot} first, is now north-west of it.
\end{solution}

\begin{solution}{exo:g7:volcanoes:9}
(a) Almost all \omterm{def:g4:volcanoes-earthquakes-first-look:volcano}{volcanoes} stand on \omterm{def:g7:plate-tectonics:plate-boundary}{plate boundaries} (87~\% in the list):
\omterm{def:g4:volcanoes-earthquakes-first-look:volcano}{volcanoes} so far from any boundary are an exception. (b) Deep in the
\omterm{def:g6:inside-the-earth:layers}{mantle} under Hawaii, a \omterm{def:g7:volcanoes:hotspot}{hotspot} stays hot for millions of years. It
partly melts the \omterm{def:g3:stones-pebbles-sand:rock}{rock} above it. The \omterm{def:g7:volcanoes:magma}{magma}, lighter than the solid \omterm{def:g3:stones-pebbles-sand:rock}{rock}
around it, rises through the \omterm{def:g7:plate-tectonics:plate}{plate} and erupts. As the \omterm{def:g7:plate-tectonics:plate}{plate} moves, new
\omterm{def:g4:volcanoes-earthquakes-first-look:volcano}{volcanoes} grow over the \omterm{def:g7:volcanoes:hotspot}{hotspot} and the old ones are carried away. (c)
Orange dots (``inside a \omterm{def:g7:plate-tectonics:plate}{plate}'').
\end{solution}

\begin{solution}{exo:g7:volcanoes:10}
The first \omterm{def:g7:volcanoes:magma}{magma} had lost most of its gas on the way up, like a bottle of
fizzy water left open or left standing: when it reached the surface,
nothing was left to blow it apart, and it oozed out. The \omterm{def:g7:volcanoes:magma}{magma} of 12 and
15~June still held its gas, like the shaken bottle: as the pressure
dropped near the surface, the gas formed bubbles all at once, which
swelled and tore the sticky \omterm{def:g7:volcanoes:magma}{magma} into ash. The first \omterm{def:g7:volcanoes:magma}{magma} is the one
that had lost its gas.
\end{solution}

\begin{solution}{exo:g7:volcanoes:11}
(a) $10 \div 80 = 0.125$~hour, and $0.125 \times 60 = 7.5$~minutes. (b)
$10 \div 130 \approx 0.077$~hour, about $0.077 \times 60 \approx
4.6$~minutes, so about 5~minutes. (c) The family needs 20~minutes, but the
flow arrives in 5 to 8~minutes: once the flow is seen, it is too late. The
family can only be safe if it leaves when the scientists see the warning
signs, before the \omterm{def:g4:volcanoes-earthquakes-first-look:eruption}{eruption}.
\end{solution}

\begin{solution}{exo:g7:volcanoes:12}
(a) $111 \times 360 = 39\,960$: about \textbf{40\,000~km}. (b)
$40\,000 \div 15 \approx 2667$: about 2\,670~km per day. (c)
$2667 \div 24 \approx 111$: about \textbf{110~km per hour}. (d) By the
wind: the fine ash has no speed of its own once the \omterm{def:g4:volcanoes-earthquakes-first-look:eruption}{eruption} has thrown
it up; it floats high in the air and the wind carries it, while it slowly
falls.
\end{solution}

\begin{solution}{pb:g7:volcanoes:1}
\textbf{1.} From 16 to 31~March: $31 - 16 = 15$~days; then 30~days of
April and 18~days of May: $15 + 30 + 18 = 63$~days.

\textbf{2.} $10\,000 \div 62 \approx 161$: more than 160 small
\omterm{def:g4:volcanoes-earthquakes-first-look:earthquake}{earthquakes} a day, on average.

\textbf{3.} The swelling showed that \omterm{def:g7:volcanoes:magma}{magma} had risen high inside the
\omterm{def:g1:landscapes-around-us:landform}{mountain} and was pushing the \omterm{def:g3:stones-pebbles-sand:rock}{rock} outwards. It also made the north side
steep and unstable, ready to slide.

\textbf{4.} A blue triangle stands where one \omterm{def:g7:plate-tectonics:plate}{plate} sinks under another. We
expect sticky, gas-rich \omterm{def:g7:volcanoes:magma}{magma}, \omterm{def:g7:volcanoes:eruption-types}{explosive eruptions} and a steep cone (a
stratovolcano) --- just what the first photograph shows.

\textbf{5.} $240 \div 60 = 4$~kilometres per minute.

\textbf{6.} The swollen side was holding the \omterm{def:g7:volcanoes:magma}{magma} in, under pressure,
like the cap of a bottle. When it slid away, the pressure dropped at
once; the gas dissolved in the \omterm{def:g7:volcanoes:magma}{magma} turned into bubbles all at once and
burst out through the open side --- sideways, where the lid had gone.

\textbf{7.} $27 \div 480 \approx 0.056$~hour, and $0.056 \times 60 \approx
3.4$~minutes. Since the blast moved at least that fast, it took at most
about 3.4~minutes.

\textbf{8.} $8 \div 130 \approx 0.062$~hour, and $0.062 \times 60 \approx
3.7$: about \textbf{4~minutes}.

\textbf{9.} $24 \div 8.8 \approx 2.7$: the column rose about 2.7 times
higher than Mount Everest.

\textbf{10.} Effusive: the sticky \omterm{def:g4:volcanoes-earthquakes-first-look:eruption}{lava} came out without exploding and piled
up over the \omterm{def:g7:volcanoes:vent}{vent}. Like the bottle that stood still, this \omterm{def:g7:volcanoes:magma}{magma} had already
lost most of its gas.

\textbf{11.} Before, the \omterm{def:g1:landscapes-around-us:landform}{mountain} was a smooth, pointed cone covered with
\omterm{def:g4:clouds-rain-snow:precipitation}{snow}, above green forests. After, its summit and north side are gone: a
huge \omterm{def:g7:volcanoes:vent}{crater} open towards us has taken their place, and the slopes and the
\omterm{def:g1:landscapes-around-us:landform}{valley} in front are grey with ash and broken \omterm{def:g3:stones-pebbles-sand:rock}{rock}.

\textbf{12.} $2950 - 2549 = 401$: the \omterm{def:g1:landscapes-around-us:landform}{mountain} lost about
\textbf{401~m} of height, about 400~metres, as told in grade~4.
\end{solution}
